\section{AMC 8 2026}
        \begin{enumerate}
        	\item (\textit{AMC 8 2026}) What is the value of the following expression?


 $1 + 2 - 3 + 4 + 5 - 6 + 7 + 8 - 9 + 10 + 11 - 12$


 $\textbf{(A)}\ 18 \qquad \textbf{(B)}\ 21 \qquad \textbf{(C)}\ 24 \qquad \textbf{(D)}\ 27 \qquad \textbf{(E)}\ 30$


 
 \textbf{Answer: }A



\textbf{Solution 1}

We can use mental calculations to get $1+2-3+4+5-6+7+8-9+10+11-12$ = $\boxed{\textbf{(A)}18}$.

 \textbf{Simple trick} To solve the expression $1 + 2 - 3 + 4 + 5 - 6 + 7 + 8 - 9 + 10 + 11 - 12$, we can define a function $f(x)$ whose derivative $f'(x)$ evaluated at $x=1$ yields the original sum. By grouping the terms into blocks of three, we have:
 \[f(x) = (x + x^2 - x^3) + (x^4 + x^5 - x^6) + (x^7 + x^8 - x^9) + (x^{10} + x^{11} - x^{12})\]
 \[f(x) = (x + x^2 - x^3)(1 + x^3 + x^6 + x^9)\]Let $u(x) = x + x^2 - x^3$ and $v(x) = 1 + x^3 + x^6 + x^9$. Using the Product Rule, $f'(x) = u'(x)v(x) + u(x)v'(x)$. Now we evaluating these components at $x=1$:

 $u(1) = 1 + 1 - 1 = 1$

 $u'(1) = 1 + 2(1) - 3(1)^2 = 0$

 $v(1) = 1 + 1 + 1 + 1 = 4$

 $v'(1) = 3(1)^2 + 6(1)^5 + 9(1)^8 = 18$

 Substituting these values into the product rule gives:
 \[f'(1) = (0 \cdot 4) + (1 \cdot 18) = \boxed{18}\]

 



\textbf{Solution 2}

$1+2-3+4+5-6+7+8-9+10+11-12$ = $\boxed{18}$

 



\textbf{Solution 3}

Note that $2-3 = 5-6 = 8-9 = 11-12 = -1$, so we have a commutive sum of $-4$. Then, the sum is just $(1+4+7+10)-4 = 1+7+10 = \boxed{18}$.

 



\textbf{Solution 4 (Break the problem into chunks!)}

Notice that we can break this up into chunks of 3 numbers. $1+2-3=0$, $4+5-6=3$, $7+8-9=6$, and $10+11-12=9$, these are all multiples of 3 so we get $3+6+9=3 \times 1 + 3\times2+3\times3= 3\times(1+2+3) = 3 \times 6 =\boxed{18}$

 



\textbf{Solution 5}

$1+2-3+4+5-6+7+8-9+10+11-12$ $=1+2+3+4+5+6+7+8+9+10+11+12-3\times2-6\times2-9\times2-12\times2$ $=\frac{12\times13}2-{3+6+9+12}\times2=$ $78-60=18$, so our answer is $\boxed{\textbf{(A)}18}$.

 



\textbf{Solution 6 (CRITICAL THINKING!!)}

Notice that $1+2-3$ and $4+5-6$ have one similarity; the latter is the former plus $3+3-3 = 3$, because $(1+3)+(2+3)-(3+3)$. This means that expression 2 is $3+3-3=3+0=3$ greater than the first, which is $3$. By trivial, $1+2-3=0.$ This means that the sum can be expressed as $1+2-3+(1+3)+(2+3)-(3+3)$. This can even be extended further to the 7+8-9 part and the 10+11-12 part, so the expression can simplify down to $0+3+3-3+6+6-6+9+9-9$, or $0+3+6+9=$ $\boxed{\text{18}}$.You can also use the pattern for the negatives. For example 1-2=-1 the negative one can really help.

 



\textbf{Solution 7 (Cancellation)}

We see that many of the numbers in the equation can be cancelled, so we take:

 $1+2-3 = 0$

 $4+5-9 = 0$

 $-6-12+8+10 = 0$

 Which leaves...

 $7+11 = \boxed{\text{18}}$

 This is a very simple method for quick, concise mental math, especially on pencil and paper tests.

 



\textbf{Solution 8 (Grouping by -1)}

You can see that from the question, you can group like: 

 $1+(2-3)+4+(5-6)+7+(8-9)+10+(11-12)$

 Which simplifies to:

 $1-1+4-1+7-1+10-1$

 And quickly, we can group it into:

 $(1-1)+(4-1)+(7-1)+(10-1)$

 and is equal to...

 $3+6+9$

 And then the answer is $\boxed{\text{18}}$

 


	\item (\textit{AMC 8 2026}) In the array shown below, $3$s are surrounded by  $2$s, which are then surrounded by a border of $1$s. What is the sum of the numbers in the array?


 
 \[\begin{array}{ccccccc} 1 & 1 & 1 & 1 & 1 & 1 & 1 \\ 1 & 2 & 2 & 2 & 2 & 2 & 1 \\ 1 & 2 & 3 & 3 & 3 & 2 & 1 \\ 1 & 2 & 2 & 2 & 2 & 2 & 1 \\ 1 & 1 & 1 & 1 & 1 & 1 & 1 \end{array}\]


 $\textbf{(A)}\ 49 \qquad \textbf{(B)}\ 51 \qquad \textbf{(C)}\ 53 \qquad \textbf{(D)}\ 55 \qquad \textbf{(E)}\ 57$


 
 \textbf{Answer: }C



\textbf{Solution 1}

Notice that if we label the rows $A_1$, $A_2$, $A_3$, ..., $A_5$ top down, then, if we denote $S(A_n)$ to be the sum of the $A_n$ row, $A_1 = A_5$, $A_2 = A_4$, and $A_3 = A_3$ through symmetry. Thus the total sum is $2A_1 + 2A_2 + A_3$.

 Then, it is obvious that the sum of $A_1 = 7$, the sum of $A_2 = 2 + 2(5) = 12$, and the sum of $A_3 = 2(3) + 3(3) = 5(3) = 15$. The answer is then $2 \cdot 7 + 2 \cdot 12 + 15 = 14 + 15 + 24 = 29 + 24 = \boxed{\textbf{(C)}\  53}$.

 



\textbf{Solution 2 (brute force)}

You can just add up all the numbers in the array, giving $\boxed{\textbf{(C)}\ 53}$. ~Baihly2024

 



\textbf{Solution 3}

We have $20$ ones, $12$ twos, and $3$ threes, so we have $20\times1+12\times2+3\times3=\boxed{\textbf{(C)}\ 53}$ as our answer.

 



\textbf{Solution 4}

Looking at the answer choices, we observe that each has a different units digit. Therefore, it is sufficient to determine the sum modulo $10$.

 There are two rows consisting of five $2$’s each (the second and fourth rows), whose total contribution is$(2+2+2+2+2)+(2+2+2+2+2)=20\equiv 0 \pmod{10}$.

 In the middle row, four of the entries sum to$2+2+3+3=10\equiv 0 \pmod{10}$,leaving one remaining $3$.

 The border of the array consists of $20$ entries equal to $1$, whose sum is$20\equiv 0 \pmod{10}$.

 Thus, the entire sum is congruent to $3 \pmod{10}$. Among the answer choices, only $\text{(C)}\ 53$ ends in $3$.

 Therefore, the sum of the numbers in the array is $\boxed{53}$.

 



\textbf{Solution 5}

The sum of the entries in this array can be found by counting each entry of the whole $5 \times 7$ array once, each entry of the inner $3 \times 5$ array a second time, and finally each entry of the innermost $1 \times 3$ array a third time.  This gives a total sum of $5 \times 7 + 3 \times 5 + 1 \times 3 = 35 + 15 + 3 = \boxed{\textbf{(C)}\  53}$

 -j314andrews

 


 



\textbf{Solution 6 (smart application of brute force)}

We can see there are $20$ '$1$'s, so we get $20 \cdot 1$ = $20$. Then, we notice there are $12$ '$2$'s, yielding a total of $12 \cdot 2$ = $24$. Finally, we notice there are $3$ '$3$'s, yielding a total of $3 \cdot 3$ = $9$. Adding all sums up, we arrive at a grand total of $20$ + $24$ + $9$ = 53. So our answer is  $\boxed{\textbf{(C)}\ 53}$.

 



\textbf{Solution 7 (Formula)}

In this solution we use a formula to find the sum of all entries in the array:

 R = Number of rows (5)

 C = Number of Columns (7)

 S = Difference between each inlaid ring of numbers ($3-2 = 1, 2-1 = 1,$ so $s = 1$)

 T = starting number: the outer ring value (1)

 $(RC - (R-2)(C-2))T + ((R-2)(C-2) - (R-4)(C-4))(T+1) + ((R-4)(C-4) - (R-6)(C-6))(T+2)$

 This of course equals $\boxed{\textbf{(C)}\ 53}$.

 You could undoubtedly simplify this using more variable, but I find this a straightforward sequence for these purposes. I understand the correlation between T, S, and the number that you subtract from the variables R and C, so this could be an opportunity for another simplification.

 I suggest you try this with more arrays to possibly come up with an easier formula.

 


	\item (\textit{AMC 8 2026}) Haruki has a piece of wire that is $24$ centimeters long. He wants to bend it to form each of the following shapes, one at a time.


 
                \begin{itemize}
                    \item A regular hexagon with side length $5$ cm.
                \end{itemize}


                
                \begin{itemize}
                    \item A square of area $36$ cm².
                \end{itemize}


                
                \begin{itemize}
                    \item A right triangle whose legs are $6$ and $8$ cm long.
                \end{itemize}


                Which of the shapes can Haruki make?


 $\textbf{(A) }\text{Triangle only} \qquad \textbf{(B) }\text{Hexagon and square only} \\  \textbf{(C) }\text{Hexagon and triangle only} \qquad \textbf{(D) }\text{Square and triangle only} \\ \textbf{(E) }\text{Hexagon, triangle, and square}$


 
 \textbf{Answer: }D



\textbf{Solution 1}

	To solve this problem we can find the perimeter of each of the given spaces and check if it is less than 24.

 The first shape is a regular hexagon with side length $5$. The hexagon has a perimeter of $6 \cdot 5 = 30$ centimeters. The second shape is a square with area $36$ or a side length of $6$. Then the perimeter of all $4$ sides is $24$ centimeters. Finally, we need to find the perimeter of a right triangle with legs $6$ and $8$ centimeters. The hypotenuse of the right triangle is $10$ so the total perimeter is $6 + 8 + 10 = 24$ centimeters.

 Thus, only the square and the triangle can be made with 24 centimeters of wire so the answer is $\boxed{\textbf {(D) } \text{Square and triangle only}}$.

 



\textbf{Solution 2 (Easy and simple, no need to know the sides of the triangle!)}

First, consider the regular hexagon. A regular hexagon with side length $5$ cm has perimeter$6\cdot 5 = 30$ cm, which exceeds the available $24$ cm of wire. Therefore, the hexagon cannot be formed, eliminating answer choices $\textbf{(B)}$, $\textbf{(C)}$, and $\textbf{(E)}$.

 Next, consider the square. A square with area $36\text{ cm}^2$ has side length $\sqrt{36}=6$ cm, so its perimeter is$4\cdot 6 = 24$ cm. Thus, the square can be formed.

 At this point, the only remaining answer choices are $\textbf{(A)}$ and $\textbf{(D)}$. But since $\textbf{(A) Triangle only}$ does not include the square and $\textbf{(D)}$ does, the correct answer must be $\textbf{(D)}$.

 Therefore, the shapes Haruki can make are the square and the triangle, and the answer is $\boxed{\textbf{(D) Square and triangle only}}$. Conveniently, the triangle is listed last, so the shapes can be evaluated in the given order.

 


	\item (\textit{AMC 8 2026}) Brynn's savings decreased by $20$\% in July, then increased by $50$\% in August. Brynn's savings are
now what percent of the original amount?


 $\textbf{(A) }80 \qquad \textbf{(B) }90 \qquad \textbf{(C) }100 \qquad \textbf{(D) }110 \qquad \textbf{(E) }120$


 
 \textbf{Answer: }E



\textbf{Solution 1}

Let $x$ be Byron's original amount. We see that Byron had $0.8x$ in July. Then, he increased his savings by HALF of how much he had as of July which is $0.4x$. Adding $0.4x + 0.8x$, we get $1.2x$, so his August savings are $\boxed{ \textbf{(E) } 120 }\%$ of $x$. 

 We see that the change is $(1-0.2)(1+0.5) = 0.8 \cdot 1.5 = 1.2$. That is an increase of $20\%$, so his August savings are $\boxed{ \textbf{(E) } 120 }\%$ of his original savings.

 



\textbf{Solution 2}

Let's assume Brynn has $\$100$ at the start. In July, he lost $20\%$, so he lost $\$20$, so he has $100-20=80$ dollars left. $50$ percent of $80$ dollars is $40$ dollars, so he gained $40$ dollars in August, so now he has $80+40=120$ dollars, which is $120\%$ of $100$ dollars, so our answer is $\boxed{ \textbf{(E) } 120 }\%$.

 We can also do it just with percentages. He has $100\%$ then loses $20\%$, meaning he now has $80\%$. Then he gains $50$ percent of our NEW amount (which is $80$ percent) and $50$ percent of $80$ is $40$ percent, so we add $40$ and $80$ to get our answer, $\boxed{ \textbf{(E) } 120 }\%$.

 You can also do it with fractions. You can say you have $\frac{100}{100}$ at the beginning, then losing $20\%$ would result in $\frac{80}{100}$. $50\%$ of $\frac{80}{100}$ is $\frac{40}{100}$. Adding $\frac{40}{100}$ to the original $\frac{80}{100}$ gives $\frac{120}{100}$, which is $\boxed{ \textbf{(E) } 120 }\%$ of $\frac{100}{100}$.

 



\textbf{Solution 3}

Let Brynn's savings be $x$. Then, in July from $x$ his savings went down $x - \frac{1}{5}x = \frac{4}{5}x$. From that amount, it increased $50\%$ in August, meaning we have $\frac{8}{10}x + \frac{4}{10}x = \frac{12}{10}x = \frac{6}{5}x$ remaining. Of the total $x$, which is $\frac{\frac{6}{5}x}{x}$, we see that the $x$ cancel out and the increase is $\frac{6}{5}$, or $\boxed{120}$\%

 


	\item (\textit{AMC 8 2026}) Casey went on a road trip that covered $100$ miles, stopping for a lunch break along the way. The trip took $3$ hours in total and her average speed while driving was $40$ miles per hour (mph). In minutes, how long was the lunch break?


 $\textbf{(A) }15 \qquad \textbf{(B) }30 \qquad \textbf{(C) }40 \qquad \textbf{(D) }45 \qquad \textbf{(E) }60$


 
 \textbf{Answer: }B



\textbf{Solution}

The average speed is $\frac{ \text{ Total distance } }{ \text{ Total time } }$, or $\frac{100}{3-t}$. Thus $\frac{100}{3-t} = 40 \Rightarrow 3-t = \frac{5}{2} \Rightarrow t = \frac{1}{2}$, which is $\boxed{ \textbf{(B) } 30}$ minutes.

 nah we can just do it like this. $100$ miles at $40$ miles per hour means we can do $1$ mile in $1.5$ minutes. so we multiply $100$ by $1.5$ to get $150$ minutes or $2$ hours and $30$ minutes. (this is how long it took her to drive, and if the total time is three hours we do $3$ minus $2.5$ to get our answer $\boxed{ \textbf{(B) } 30}$ minutes.

 thanks pookies~pereena

 



\textbf{Solution 2}

Average speed is total distance over total time. Thus, the total time spent driving is $t$ and the total distance is $100$, the average speed is $100/t = 40$. Solving, we have that $t = 100/40 = 10/4 = 5/2$ hours, or $2$ hours and $30$ minutes. She drove for $3$ hours, and thus her lunch break was $\boxed{30}$ minutes.

 



\textbf{Solution 3}

If the trip took $3$ hours in total and her average speed was $40$ miles per hour, we know she would have traveled $120$ miles. Since she only traveled $100$ miles with a lunch break, she would have traveled $20$ miles if she didn't have a lunch break. Since her average speed was $40$ miles per hour, the lunch break took $60*\frac{20}{40}$ minutes. So the answer is $\boxed{30}$ minutes.

 



\textbf{Solution 4}

Since Casey drove $100$ miles at a constant rate of $40$ mph, we can say it took her $100/40 = 2.5$ hours to drive to his/her destination WITHOUT a break. The problem mentioned the trip took $3$ hours, so $3 - 2.5$ hours yields a total of $0.5$ hours or $30$ minutes, so the answer is $\boxed{ \textbf{(B) } 30}$.

 



\textbf{Solution 5}

Using dimensional analysis, her time spent driving is
 \[100 \text{ miles}\cdot \frac{ 1 \text{ hour} }{ 40 \text{ miles} } = 2.5 \text{ hours}.\]Therefore, she spent $3 - 2.5 = 0.5 \text{ hours}$ during her lunch break or
 \[0.5 \text{ hours} \cdot \frac{60 \text{ minutes}}{1 \text{ hour}} = \boxed{ \textbf{(B) } 30} \text{ minutes}.\]

 


	\item (\textit{AMC 8 2026}) Peter lives near a rectangular field that is filled with blackberry bushes. The field is $10$ meters long and $8$ meters wide, and Peter can reach any blackberries that are within $1$ meter of an edge of the field. The portion of the field he can reach is shaded in the figure below. What fraction of the area of the field can Peter reach?


 
                    \begin{figure}[H]
                        \centering
                        \includegraphics[width=0.5\textwidth]{Problem/AMC_8_AMC_2026_6_Problem_Figure_1.png}
                    \end{figure}
                    $\textbf{(A) } \frac{1}{6}\qquad\textbf{(B) } \frac{1}{4}\qquad\textbf{(C) } \frac{1}{3}\qquad\textbf{(D) } \frac{3}{8}\qquad\textbf{(E) } \frac{2}{5}\qquad$


 
 \textbf{Answer: }E



\textbf{Solution}

We can calculate this reachable area to be $\text{Total area} - \text{Unreachable area}$. The total area is $8 \cdot 10 = 80$, and the unreachable area is $(8-2)(10-2) = 6\cdot 8 = 48$. Thus, our answer is $\frac{80-48}{80} = \frac{32}{80} = \boxed{ \textbf{(E) } \frac{2}{5} }$.

 



\textbf{Solution 2}

Each "region" is one unit from an edge. Thus, let $l$ be the length and $w$ be the width of the unshaded region. Then, it is obvious that $l = 10 - 2 = 8$ and $w = 8 - 2 = 6$, and the unshaded regions dimensions are $8 \times 6$.

 Then, the area of the shaded region is the area of the total region subtracted by the area of the unshaded region, or $80 - 48 = 32$ $\text{un}^2$. Then, the ratio is simply $32/80$, or $\boxed{\frac{2}{5}}$.

 



\textbf{Solution 3}

By rearranging the diagram, the unreachable (white) region can be moved to the lower-left corner without changing the shaded area. The reachable region then forms an L-shaped figure composed of two rectangles.

 One rectangle has dimensions $8\times 2$ and the other has dimensions $10\times 2$. These rectangles overlap in a $2\times 2$ square, which must be subtracted once.

 Thus, the area of the shaded region is$(8\cdot 2)+(10\cdot 2)-(2\cdot 2)=2(10+8-2)$.

 The total area of the field is $10\cdot 8=80$. Dividing, the fraction of the field Peter can reach is$\dfrac{2(10+8-2)}{80}=\dfrac{10+8-2}{40}=\dfrac{16}{40}=\boxed{\textbf{(E)}\ \dfrac{2}{5}}$.

 



\textbf{Solution 4}

We can first start off by finding the unreachable area. The problem tells us that the distance between the reachable and unreachable area is $1$, and since there are $2$, we get the dimensions of the unreachable to be $8 \cdot 6$, because $8 = 10 - 2$, and $6 = 8 - 2.$. Therefore, the area of the unreachable is $48$. However, the problem asks us to find the fraction of area he can reach, so we need to get area of the entire field and subtract by $48$.Since the entire field is $8 \cdot 10 = 80$,  we get $80 - 48 = 32$. So the area of the part Peter can reach is $32$. Finally, to convert this to a fraction of the entire field, we can do $\frac{32}{80}$, so our answer is $\frac{32}{80} = \boxed{ \textbf{(E) } \frac{2}{5} }$.

 To add on, its just like calculating the total area- the unwanted recangular area.

 


	\item (\textit{AMC 8 2026}) Mika would like to estimate how far she can ride a new model of electric bike on a fully charged battery. She completed two trips totaling $40$ miles. The first trip used $\frac{1}{2}$ of the total battery power, while the second trip used $\frac{3}{10}$ of the total battery power. How many miles can this electric bike go on a fully charged battery?


 $\textbf{(A) } 45\qquad\textbf{(B) } 48\qquad\textbf{(C) } 50\qquad\textbf{(D) } 52\qquad\textbf{(E) } 55\qquad$


 
 \textbf{Answer: }C



\textbf{Solution}

We can see that the total battery spent is $\frac{1}{2} + \frac{3}{10} = \frac{5+3}{10} = \frac{4}{5}$. We can thus write the proportion $\frac{40}{\frac{4}{5}} = \frac{x}{1}$. Since $\frac{40}{\frac{4}{5}} = 40\cdot \frac{5}{4} = \boxed{50}$

 



\textbf{Solution 2}

The fully charged battery has a total battery power $P$. There is a bijective analysis that if one uses a set speed for both ways in a directional plane, the amount $P$ is the farthest they can travel. Then, after the first trip, $\frac{1}{2}P$ is used. Then, after the second, $\frac{3}{10}P$ is used. Thus, the total distance traveled (40 miles) corresponds to the total fraction of battery used:
 \[\frac{1}{2}P + \frac{3}{10}P = 40\]To clear the denominators, we multiply by 10:
 \[5P + 3P = 400\]
 \[8P = 400\]
 \[P = 50\]The answer is clearly $\boxed{50}$.

 



\textbf{Solution 3}

Start with combining the battery fractions listed in the problem, and we compare that to the number of miles she rode.

 $\frac{1}{2} + \frac{3}{10} = \frac{8}{10}$

 So $40$ miles is $\frac{8}{10}$ of her battery power.

 
 \[\frac{8}{10}x = 40\]

 Multiply both sides by $\frac{10}{8}$:

 
 \[\frac{10}{8} \cdot \frac{8}{10}x = 40 \cdot \frac{10}{8}\]

 
 \[x = \frac{400}{8}\]

 The answer is $\boxed{50}$.

 


	\item (\textit{AMC 8 2026}) A poll asked a number of people if they liked solving mathematics problems. Exactly $74$\% answered “yes.” What is the fewest possible number of people who could have been asked the question?


 $\textbf{(A)}\ 10 \qquad \textbf{(B)}\ 20 \qquad \textbf{(C)}\ 25 \qquad \textbf{(D)}\ 50 \qquad \textbf{(E)}\ 100$


 
 \textbf{Answer: }D



\textbf{Solution 1}

Well, we use our critical thinking skills to know that you can’t have .$5$ or .$7$ people (or any fraction for that matter). So, we will be like me while taking this test and do each answer individually (starting from smallest) until one gives us a whole number. which in this case that is $\boxed{ \textbf{(D) } 50}$. (although this is quite slow so I’d use the other methods posted)~pereena come on let's think positively, ok?

 



\textbf{Solution 2}

Let $n$ be the number of people. Then, $\frac{74}{100} n$ like to do math (I think I am one of them). $\frac{74}{100}$ simplifies down to $\frac{37}{50}$, so the remaining people who don't like math (haters just hate I guess) would also fall under the $50$ denominator category. The answer is $\boxed{50}$

 



\textbf{Solution 3}

We start by trying every answer from smallest to greatest. We have that $10$ is the smallest choice, so we try this. If we polled $10$ people, then $7.6$ must have said yes, which is clearly impossible, since fractional people can't exist so we discard this answer. If we continue in this matter, we can find that the answer is $\boxed{ \textbf{(D) }50}$. 

 



\textbf{Solution 4}

If $10$ people are polled, then $74$ must be divisible by $10$, but it's not, so it is not $10$ people. If 20 people are polled, then $74$ must be divisible by $5$, but it's not, so it is not $20$ people. If $25$ people are polled, then $74$ must be divisible by $4$, but it's not, so it is not $25$ people. If $50$ people are polled, then $76$ must be divisible by $2$, which it is, so our answer is $\boxed{ \textbf{(D) }50}$. 

 



\textbf{Solution 5}

Our percent is 74Assuming the people are in whole numbers\ldots

 If it is A) 10, then the percent will have to be a multiple of $\frac{100}{10}$ = 10

 If it is B) 20, then the percent will have to be a multiple of $\frac {100}{20}$ = 5

 If it is C) 25, then the percent will have to be a multiple of $\frac {100}{25}$ = 4

 It it is D) 50, then the percent will have to be a multiple of $\frac {100}{50}$ = 2

 If it is E) 100, then the percent will have to be a multiple of $\frac {100}{100}$ = 1

 In this case, 74 only fits D and E, and because D is less, $\boxed{ \textbf{(D) }50}$ is the correct answer.

 



\textbf{Solution 6}

We can see that $74 \% = \frac{74}{100}$ by definition. This fraction can be simplified to $\frac{37}{50}$, meaning $37$ out of $50$ people said "Yes". Since $\gcd(37,50) = 1$, if this fraction was divided any further, we would have fractional numerator and denominator, which is clearly impossible. Therefore, the minimum number of people surveyed was $\boxed{ \textbf{(D) } 50}$.~Ingeniouspeacock36

 


	\item (\textit{AMC 8 2026}) What is the value of this expression?


 
 \[\frac{\sqrt{16\sqrt{81}}}{\sqrt{81\sqrt{16}}}\]


 $\textbf{(A) } \frac{4}{9}\qquad\textbf{(B) } \frac{2}{3}\qquad\textbf{(C) } 1\qquad\textbf{(D) } \frac{3}{2}\qquad\textbf{(E) } \frac{9}{4}$


 
 \textbf{Answer: }B



\textbf{Solution 1}

The inner square root cannot be negative. If the outer square root is negative, the denominator is also negative, resulting in a positive fraction, an identity. Thus, we only consider positive squares, in which the numerator simplifies to $\sqrt{16 \cdot 9} = \sqrt{16}\sqrt{9} = 4 \cdot 3 = 12$, and the denominator simplifies to $\sqrt{81 \cdot 4} = \sqrt{81}\sqrt{4} = 9 \cdot 2 = 18$. The answer is $\frac{12}{18} = \frac{4}{6} = \boxed{\frac{2}{3}}$.

 (Other solutions weren't explicit with logic).

 



\textbf{Solution 2}

Since $\sqrt{81} = 9$ and $\sqrt{16} = 4$, we can simplify this expression to $\frac{\sqrt{16\cdot 9}}{\sqrt{81\cdot 4}} = \frac{\sqrt{144}}{\sqrt{324}}$, which can be calculated to be $\frac{12}{18} = \boxed{ \textbf{(B) } \frac{2}{3}}$.

 



\textbf{Solution 3}

$\frac{\sqrt{16\sqrt{81}}}{\sqrt{81\sqrt{16}}}=\frac{4\sqrt{\sqrt{81}}}{9\sqrt{\sqrt{16}}}=\frac{4\times3}{9\times2}=\frac{12}{18}=\frac23$, so our answer is $\boxed{ \textbf{(B) } \frac{2}{3}}$.

 



\textbf{Solution 4}

$\frac{\sqrt{16\sqrt{81}}}{\sqrt{81\sqrt{16}}}=\frac{\sqrt{16\cdot9}}{\sqrt{81\cdot4}}$ Notice that $16=4^2$, $9=3^2$, $81=9^2$, and $4=2^2$. We also know that $x^2 y^2 = (xy)^2$. So this equals $\frac{\sqrt{4^2\cdot3^2}}{\sqrt{9^2\cdot2^2}}=\frac{\sqrt{(4\cdot3)^2}}{\sqrt{(9\cdot2)^2}}=\frac{\sqrt{12^2}}{\sqrt{18^2}}=\frac{12}{18}=\frac{2}{3}$, so our answer is $\boxed{ \textbf{(B) } \frac{2}{3}}$.

 



\textbf{Solution 5}

Since $\sqrt[3]{81} = 3$, and $\sqrt{16} = 4$, the numerator of $\frac{\sqrt{16\sqrt{81}}}{\sqrt{81\sqrt{16}}}$ becomes $4 \cdot 3 = 12$. To find the denominator's value, we know that $\sqrt[3]{16} = 2$, and that $\sqrt{81} = 9$, the denominator of $\frac{\sqrt{16\sqrt{81}}}{\sqrt{81\sqrt{16}}}$ yields $9 \cdot 2 = 18$. We are left with $\frac{12}{18}$. After further simplification, our answer is $\boxed{ \textbf{(B) } \frac{2}{3}}$.~Ingeniouspeacock36

 


	\item (\textit{AMC 8 2026}) Five runners completed the grueling Xmarathon: Luke, Melina, Nico, Olympia, and Pedro.


 
                \begin{itemize}
                    \item Nico finished $11$ minutes behind Pedro.
                \end{itemize}


                
                \begin{itemize}
                    \item Olympia finished $2$ minutes ahead of Melina, but $3$ minutes behind Pedro.
                \end{itemize}


                
                \begin{itemize}
                    \item Olympia finished $6$ minutes ahead of Luke.
                \end{itemize}


                Which runner finished fourth?


 $\textbf{(A) }\text{Luke} \qquad \textbf{(B) }\text{Melina} \qquad \textbf{(C) }\text{Nico} \qquad \textbf{(D) }\text{Olympia}\qquad \textbf{(E) }\text{Pedro}$


 
 \textbf{Answer: }A



\textbf{Solution 1}

Let the respective places be $L$, $M$, $N$, $O$, and $P$ respectively. Then, $P = 11+N$, $O-2 = M$, $O+3 = P$, and $O = 6+L$. Then, substitution for $P$ gives us $O+3 = N+11$, in which $O=N+8$. Then, $N+6 = M$, $N+8=6+L \Rightarrow N+2=L$. So $N$ finished m minutes behind everyone (that was ok), and thus was last.

 Now, $O=6+L$, in which 

 


	\item (\textit{AMC 8 2026}) Squares of side length $1$, $1$, $2$, $3$, and $5$ are arranged to form the rectangle shown below. A curve is drawn by inscribing a quarter circle in each square and joining the quarter circles in order, from shortest to longest. What is the length of the curve?


 
                    \begin{figure}[H]
                        \centering
                        \includegraphics[width=0.5\textwidth]{Problem/AMC_8_AMC_2026_11_Problem_Figure_1.png}
                    \end{figure}
                    $\textbf{(A) } 4\pi\qquad\textbf{(B) } 6\pi\qquad\textbf{(C) } \frac{13}{2}\pi\qquad\textbf{(D) } 8\pi\qquad\textbf{(E) } 13\pi$


 
 \textbf{Answer: }B



\textbf{Solution 2 (No Fibonacci)}

For those who don't know Fibonacci's Spiral, One could see that between the two unit squares, the curvature is $\frac{1}{2} \cdot 2 \pi (1) = \pi$. For the unit 2 square, the curve is $\frac{1}{4} \cdot 2 \pi (2) = \pi$. For the unit three square, is is $\frac{1}{4} \cdot 2 \pi (3) = \frac{3}{2} \pi$. For the unit 5 square, we have $\frac{1}{4} \cdot 2 \pi \cdot 5 = \frac{10}{4} \pi = \frac{5}{2} \pi$. The length of the curve is $2\pi + \frac{8}{2} \pi = 2\pi + 4\pi = \boxed{6\pi}$.

 



\textbf{Solution 3 (Fibbonacci)}

For those who know Fibonacci's Spiral, the length of the curve $L_n = \frac{\pi}{2} \sum_{k=1}^n F_k$ where $F_k$ is the $k$th Fibonacci Number. Here, $n$ is the number of squares the curve traverses. Then, $L_5 = \frac{\pi}{2} \cdot (12) = \boxed{6 \pi}$.

 



\textbf{Solution 4 (Straightforward)}

Note that the circumference of a quarter arc is $\frac{1}{4} \cdot 2 \cdot \pi \cdot r$Which is the same thing as $\frac {1}{2} \cdot \pi \cdot r$

 The first part of the spiral has a quarter circle with the radius of 5 and so the formula would be 
 \[\frac{1}{2} \cdot \pi \cdot 5 = \frac{5}{2}\pi\]

 The second part of the spiral has a quarter circle with the radius of 3 and so the formula is 
 \[\frac{1}{2} \cdot \pi \cdot 3 = \frac{3}{2}\pi\]

 The third part of the spiral has a quarter circle with the radius of 2 and so the formula is 
 \[\frac{1}{2} \cdot \pi\cdot2 = \pi\]

 The fourth and fifth part can be viewed as 2 quarter circles with radius 1, so they are essentially one semicircle with radius 1. The formula then would be $\pi \cdot \frac{1}{2} \cdot 2 \cdot r = \pi$ Therefore, it is just $\pi$. 

 Adding them all up, we get 
 \[\frac{5}{2}\pi + \frac{3}{2}\pi + 1\pi +1\pi = \frac{8}{2}\pi + 2\pi = 6\pi\]Therefore our final answer is  
 \[\boxed{6\pi}.\]

 -Razorpro10

 


	\item (\textit{AMC 8 2026}) In the figure below, each circle will be filled with a digit from $1$ to $6$. Each digit must appear exactly once.
The sum of the digits in neighboring circles is shown in the box between them.
What digit must be placed in the top circle?


 
                    \begin{figure}[H]
                        \centering
                        \includegraphics[width=0.5\textwidth]{Problem/AMC_8_AMC_2026_12_Problem_Figure_1.png}
                    \end{figure}
                    


 $\textbf{(A) } 2\qquad\textbf{(B) } 3\qquad\textbf{(C) } 4\qquad\textbf{(D) } 5\qquad\textbf{(E) } \text{it is impossible to fill the circles}$


 
 \textbf{Answer: }D



\textbf{Solution}

Let $O$ denote an odd number and $E$ denote an even. Then, we know that $O+E = O$, and thus the two neighboring circles that form sums 9 and 5 must be one of each odd and even. To obtain 5, we have $1+4$ or $2+3$. To obtain 9, we have $3+6$, $4+5$. Consider each case separately.

 Case 1: If $1+4$ makes the sum of 5, then the two numbers that make 9 must be $3+6$. Then, 3 is forced on top, and we require another 3. However, the problem doesn't allow this, so this is impossible.

 Case 2: $2+3$ is on 5, meaning that $4+5$ is on 9. Then, we have two subcases.

 Case A: 4 is on top. Then, we require 2 to complete the sum of 6, which is not possible as $2+3$ is on 5.

 Case B: 5 is on top, meaning that we require 1 right below, 3 on the bottom right, 2 to complete the sum of 5, and the rest are uniquely determined.

 The number is $\boxed{5}$.

 



\textbf{Solution 2}

In the bottom-left corner, we have 2 numbers that add up to 10 filled in from the digits from $1$ to $6$. There are two ways you can do this, $5$ and $5$ or $6$ and $4$. Since we can only use each integer once, the first case doesn't work, so we must have the second case. Now we need to figure out the ordering. If we have $4$ by the $8$, the other circle must be $4$, but again, we can only use each positive integer once, so the $6$ must be by the $8$, and the $4$ must be by the $10$ and the $9$. Notice that we can already find the number that needs to go into the circle we want, which is $\boxed{\textbf{(D) }5}$ since $9-4=5$. ~Baihly2024

 



\textbf{Solution 3}

We do trial and error, and because it is a shorter problem it doesn’t take long. we can find that the only possibility is $\boxed{\textbf{(D) }5}$。

 



\textbf{Solution 4 (Prob the fastest)}

Let's set space 1 to the top, space 2 to the middle left, space 3 to the middle right, space 4 to the bottom left, space 5  to the bottom middle, and space 6 to the bottom right. Note that I am only doing this to better explain this to you, the numbering of spaces is just a waste of time on the AMC 8.

 Look at space 3 and space 6. The numbers must sum to 4. The two numbers cannot be 2 and 2, as each digit must repeat only once, therefore the only number pair that satisfies this is 1 and 3. Now we must figure out where to place these numbers. We can now look at space 1 and space 3. If we say that space 3 is equal to 1, then space 1 must be equal to 5. That is one possible solution. Now, if we made space 3 equal to 3 instead, space 1 must be equal to 3 as well, which contradicts the problem. Therefore, the only possible solution is 5. Now, we must check and solve the rest of the triangle to make sure that there is a solution, and the answer is not E. We can now find that the final answer is $\boxed{\textbf{(D) }5}$.

 


	\item (\textit{AMC 8 2026}) The figure below shows a tiling of 1×1 unit squares. Each row of unit squares is shifted horizontally by half a unit relative to the row above it. A shaded square is drawn on top of the tiling. Each vertex of the shaded square is a vertex of one of the unit squares. In square units, what is the area of the shaded square?


 
                    \begin{figure}[H]
                        \centering
                        \includegraphics[width=0.5\textwidth]{Problem/AMC_8_AMC_2026_13_Problem_Figure_1.png}
                    \end{figure}
                    


 $\textbf{(A) } 10\qquad\textbf{(B) } \frac{21}{2}\qquad\textbf{(C) } \frac{32}{3}\qquad\textbf{(D) } 11\qquad\textbf{(E) } \frac{34}{3}$


 
 \textbf{Answer: }A



\textbf{Solution}

Label the vertices $A$, $B$, $C$, $D$, in clockwise order. Then, we use coordinate geometry to easily solve this. Let point $A$ be on $(0, 0)$. Then, point $B$ is 1 unit up and 3 right of $(0, 0)$, or $(3, 1)$. The figure is said to be a square. and so we just calculate the distance from $(0, 0)$ to $(1, 3)$. Using the distance formula, we have $\sqrt{(1-0)^2 + (3-0)^2} = \sqrt{(1+9)} = \sqrt{10}$. The area is just this value squared, or $\boxed{10}$.

 



\textbf{Solution 2}

Since the side length of a square is 1 unit, we can use Pythagoras on a side length (note that we only need to find one side since the shape is a square). If we do it on the left side, we get $\sqrt{3^2+(0.5+0.5)^2}=\sqrt{10}$. Therefore, the final area is the side length squared, which is $(\sqrt{10})^2=\boxed{\textbf{(A) }10}$. ~Baihly2024

 



\textbf{Solution 3}

Notice that the top-left side of our square is the hypotenuse of a right triangle with legs of length $1$cm and $3$cm. Using the Pythagorean Theorem, the side length of our square is $\sqrt{1^2+3^2}=\sqrt{10}$cm, so the area is $\sqrt{10}^2=10$cm$^2$, so we obtain $\boxed{\textbf{(A) }10}$.

 



\textbf{Solution 4(Quick And Easy)}

Notice that we can construct a 4x4 square that passes tangent to each corner. The area of this square is 16. This square forms 4 congruent right triangles, with a base of 3 and height of 1. The area of one of these triangles is 1.5, and you can multiply that by 4 to get the total area of these triangles, 6. Subtract that from the area of the constructed square to get $\boxed{\textbf{(A) }10}$. as your answer.

 It would look something like  
                    \begin{figure}[H]
                        \centering
                        \includegraphics[width=0.5\textwidth]{Solution/AMC_8_AMC_2026_13_Solution_Figure_1.png}
                    \end{figure}
                    

 image by -N.

 



\textbf{Solution 5}

Similar to Solution 4, we construct one right triangle outside the square. We can see that it has side lengths 3 and 1. Using the Pythagorean Theorem we find that the hypotenuse, which is also a side of the original square, is $\sqrt{10}$. Now, we can find the area of the square by squaring that, which yields $\boxed{\textbf{(A) }10}$.

 



\textbf{Solution 6 (Without Pythagoras or Distance Formula)}


                \begin{itemize}
                    \item Note: this solution is unreliable and only use it if you have no other methods to solve.
                \end{itemize}


                We can count the number of small squares inside the large square and add $1$ to our running count if the square is more than half shaded, add $\frac{1}{2}$ if it is exactly (or almost exactly) $\frac{1}{2}$ shaded, and add $0$ if it is less than half shaded. Counting the squares, we have $5$ fully shaded, $4$ more than half, and $2$ exactly half. $5+4+2\times{\frac{1}{2}}=\boxed{\textbf{(A)} 10}$

 


	\item (\textit{AMC 8 2026}) Jami picked three equally spaced integer numbers on the number line. The sum of the first and the second numbers is $40$, while the sum of the second and third numbers is $60$. What is the sum of all three numbers?


 $\textbf{(A)}\ 70 \qquad \textbf{(B)}\ 75 \qquad \textbf{(C)}\ 80 \qquad \textbf{(D)}\ 85 \qquad \textbf{(E)}\ 90$


 
 \textbf{Answer: }B



\textbf{Solution}

Since the integers are equally spaced, they form an arithmetic sequence, so we can label the middle term $a$, the smallest one $a-d$ and the largest one $a+d$. We are given that $(a-d)+a=2a-d=40$ and $a+(a+d)=2a+d=60$. Solving, $(2a-d)+(2a+d)=4a=100$ so $a=25$. The problem asks for the sum of all 3 numbers, which is $(a-d)+a+(a+d)=3a=3 \cdot 25 = \boxed{\textbf{(B)} 75}$.

 



\textbf{Solution 2}

The numbers $a$, $b$, and $c$ are equally spaced by a distance $k$ on the number line. Then, we have $b = a+k$ and $c = a+2k$. Then, the sum of the first and second numbers is $a+a+k = 40 = 2a+k = 40$, and the sum of the second and third numbers is $a+k+a+2k = 2a+3k = 60$. Then, solving we obtain $-2k = -20$, $k = 10$, and thus $2a+10 = 40$, $a= 15$. Then, $b=25$, and $c=35$. The sum of all three is $15+25+35 = 70+5 = \boxed{\textbf{(B)} 75}$.

 



\textbf{Solution 3}

The difference of the sum of the first two integers and the sum of the last two integers differs by $20$. Hence, the common difference of the arithmetic sequence must be $10$. The only way for the sum of two integers in an arithmetic sequence with common difference 10 to have a sum ending in $0$ is for all the integers to end in either $0$ or $5.$

 Testing $5$ and $15$ and $10$ and $20$ does not meet the required sum for the first two integers in the sequence. However, the sum of $15$ and $25$ is equivalent to $40,$ and adding $35 (25+10)$ to $25$ matches up with the sum of the last two integers, $60.$ Consequently, the arithmetic sequence is $15, 25, 35,$ with a common difference of $10.$ Adding up all the terms in the sequence gives the answer, $15+25+35=\boxed{\textbf{(B)} 75}$

 



\textbf{Solution 4, (the most complicated, just math)}


 \[x + y = 40\] (equation 1)
 \[y + z = 60\] (equation 2)
 \[x = a, y = a + d, z = a + 2d\] 
 \[x + x + d = 40\] (equation 1)
 \[x + d + x + 2d = 40\] (equation 2)
 \[(2x + 3d) - (2x + d) = 20\] (eqs. 1 and 2 subtracted)
 \[2d = 20, d = 10\]
 \[2a + 10 = 40, a = 15\] (substitute)
 \[a + a + d + a + 2d = 15 + 15 + 10 + 15 + 20 = \boxed{75}\] (plug in and substitute)

 


	\item (\textit{AMC 8 2026}) Elijah has a large collection of identical wooden cubes which are white on 
$4$ faces and gray on 
$2$ faces that share an edge. He glues some cubes together face-to-face. The figure below shows 
$2$ cubes being glued together, leaving 
$3$ gray faces visible. What is the fewest number of cubes that he could glue together to ensure that no gray faces are visible, no matter how he rotates the figure?


 
                    \begin{figure}[H]
                        \centering
                        \includegraphics[width=0.5\textwidth]{Problem/AMC_8_AMC_2026_15_Problem_Figure_1.png}
                    \end{figure}
                    


 $\textbf{(A)}\ 4 \qquad\textbf{(B)}\ 6 \qquad\textbf{(C)}\ 8 \qquad\textbf{(D)}\ 9 \qquad\textbf{(E)}\ 27$


 
 \textbf{Answer: }A



\textbf{Solution 1 (2D Plane)}

Create a square in which two adjecent sides are colored red. Now, connect each of $n$ squares such that each red side is tangent to another. What is the minimum number required?

 By solving this bijection, we solve the problem above. Through experimentation, each square, assuming the colored sides are the rightmost and bottom of a square, must have another square on the side and below. Thus we obtain 3 squares.

 However, the top right square must have one more square below it, giving us a final count of $3+1 = \boxed{4}$ minimum cubes.

 



\textbf{Solution 2 (omg so easy)}

We can have $4$ cubes with all their gray faces facing each other, and that will make no gray faces visible. Therefore, our answer is $\boxed{\textbf{(A)}\ 4}$. Note that this is also the minimum answer choice.

 



\textbf{Solution 3}

Every time we add a new cube, there are $2$ colored faces that are added, but only $1$ colored face on $2$ cubes that are covered. If we start with $1$ cube, we see adding another cube colored-on-colored can cause both remaining colored faces to face in the same direction. So if we have $2$ $2$ x $1$ cube formations, by aligning all of the colored faces onto one plane we get the minimum of $\boxed{4}$. That is pretty nice beccause we won't have any doubts about whether it could be lower.

 


	\item (\textit{AMC 8 2026}) Consider all positive four-digit integers consisting of only even digits. What fraction of these integers are divisible by 4?


 $\textbf{(A) } \frac{1}{4} \qquad \textbf{(B) } \frac{2}{5} \qquad \textbf{(C) } \frac{1}{2} \qquad \textbf{(D) } \frac{3}{5} \qquad \textbf{(E) } \frac{3}{4}$


 
 \textbf{Answer: }D



\textbf{Solution 1}

The divisibility rule for $4$ is that the last 2 digits must be a multiple of $4$. Because of this, we can "discard" the first two digits, as the last two digits are the only ones that matter (this doesn't actually change the answer, since for every possible case that the first two digits can be, there are the same number of cases the last two digits can be, so the ratio of cases that work to the cases that don't work is still the same). We can also notice that the tens digit also doesn't matter, since every single digit in the number must be even. Therefore, we only need to care about the ones digit. We have the even digits $0$, $2$, $4$, $6$, and $8$, and there are 3 digits that are divisible by $4$, namely $0$, $4$, and $8$, out of $5$ total digits, so we have the final fraction $\boxed{\textbf{(D) }\frac{3}{5}}$. ~Baihly2024

 



\textbf{Solution 2}

A number is divisible by 4 if and only if the last two digits are divisible by 4. Then, the third and fourth digits must be divisible by 4. So, those possibilities are $\{(00), (04), (08), (12), (16), (20),\dots,(96)\}$. There are 25 cases in this, however, some of these are odd! Luckily, each odd tens digit would have an even units digit of either 2 or 6 on the third or fourth multiples, meaning each sequence of 5 between 4 and 20 actually contributes 3 to the count, giving us $5(3) - 1 + 1 = 15$ actual cases.

 Then, the remaining two numbers are to be consisting of the digits $\{0,2,4,6,8\}$. The first digit has 4 choices, and the next has 5, giving a total of 20 unique combinations. These 20 combinations can go with any of the other 24 combinations, giving 300 total combinations. There are a total of $4 \times 5 \times 5 \times 5 = 125 \times 4 = 500$ four digit integers whose digits are all even. The fraction is $300/500 = \boxed{3/5}$

 


	\item (\textit{AMC 8 2026}) Four students are seated in a row. They chat with the people sitting next to them, then rearrange themselves so that they are no longer seated next to any of the same people. How many rearrangements are possible?


 
                    \begin{figure}[H]
                        \centering
                        \includegraphics[width=0.5\textwidth]{Problem/AMC_8_AMC_2026_17_Problem_Figure_1.png}
                    \end{figure}
                    


 $\textbf{(A)}\ 2 \qquad\textbf{(B)}\ 4 \qquad\textbf{(C)}\ 9 \qquad\textbf{(D)}\ 12 \qquad\textbf{(E)}\ 24$


 
 \textbf{Answer: }A



\textbf{Solution 1 (AMC 10 Bijection)}

Label the human beings A, B, C, and D in that order from left to right. It doesn't matter the original configuration. Now, the problem is:

 How many four letter strings ABCD exist such that A is not next to B, B not next to C, and C is not next to D.

 We can count the number of strings manually (or just do casework), to obtain $\boxed{2}$.

 This is incredibly similar to one AMC 10 problem, that asked how many four letter strings abcd have no two consecutive letters in the alphabet next to one another.

 



\textbf{Solution 2}

Let's say the people are $A$, $B$, $C$, and $D$, and they seated in the order $ABCD$ at first. $A$ cannot be next to $B$, $B$ cannot be next to $C$, and $C$ cannot be next to $D$, so our only possibilities are $CADB$ and $BDAC$ as orders, so our answer is $\boxed{\mathbf{(A) 2}}$.

 



\textbf{Solution 3}

Since $B$ sits next to both $A$ and $C$, $B$ can only sit next to $D$. The only way this is possible is if $B$ sits at either the front or the back. We can apply this same logic to $C.$ There are now two cases: $BDAC$ or $CADB,$ which is the former with its order reversed. Therefore, our answer is $\boxed{\textbf {(A)} 2}$.

 


	\item (\textit{AMC 8 2026}) In how many ways can $60$ be written as the sum of two or more consecutive odd positive integers that are arranged in increasing order?


 $\textbf{(A) } 1\qquad\textbf{(B) } 2\qquad\textbf{(C) } 3\qquad\textbf{(D) } 4\qquad\textbf{(E) } 5$


 
 \textbf{Answer: }B



\textbf{Solution 1 (Casework)}

Let $n$ be the first number in such a series, and let $k$ be the number of terms in the series.  Then the sum of the series is $n + (n+2) + ... + (n+2(k-1)) = kn + 2\left(\frac{(k-1)k}{2}\right) = kn + k(k-1)$.  The goal is thus to find all solutions to the equation $kn + k(k-1) = 60$, where $n$ is an odd positive integer and $k > 2$ is an integer.

 Consider the following cases:

 Case 1: $k=2$.Then $2n+2 = 60$, which yields a valid solution in $n = 29$ and the series $29 + 31$.

 Case 2: $k=3$.Then $3n + 6 = 60$, which yields an invalid solution in $n = 18$.

 Case 3: $k=4$.Then $4n+12 = 60$, which yields an invalid solution in $n = 12$.

 Case 4: $k=5$.Then $5n+20 = 60$, which yields an invalid solution in $n = 8$.

 Case 5: $k=6$.Then $6n+30 = 60$, which yields a valid solution in $n = 5$ and the series $5 + 7 + 9 + 11 + 13 + 15$.

 Case 6: $k=7$.Then $7n + 42 = 60$, which yields an invalid solution in $n = \frac{18}{7}$.

 Case 7: $k = 8$.Then $8n + 56 = 60$, which yields an invalid solution in $n = \frac{1}{2}$.

 Case 8: $k \geq 9$.Then $nk + k(k-1) > k(k-1) > 60$, so there will not be any more valid solutions.

 Therefore, there are $\boxed{\textbf{(B) }2}$ such series.

 



\textbf{Solution 2}

We can write the sum of $n$ integers as 
 \[a + (a + 2) + (a + 4) + \dots + (a + 2(n-1)).\] This sum can be simplified algebraically to $an + 2 \cdot \frac{n(n-1)}{2} = an + n(n-1) = n(a + n - 1)$. Thus, 
 \[n(a + n - 1) = 60.\]

 We are given that $a$ is odd, along with all the other terms. If $n$ was odd as well, we would have the sum be odd as well, and $60$ is clearly not odd. Thus, we know $n$ is even.

 Trying different values of $n$ from $2$ onwards, we find the following valid ordered pairs for $(a, n)$: $(29, 2), (5, 6)$. Therefore, there are only $\boxed{\textbf{(B) } 2}$ solutions that meet the conditions of the problem.

 Note: The solutions we found here can be written as $29 + 31$ and $5 + 7 + 9 + 11 + 13 + 15$.

 



\textbf{Solution 3 (Manual Ordered Pairing [MOP])}

A sum of consecutive odd positive integers can also be represented as an arithmetic sequence, in this case, with common difference 2 and the first term being odd. The sum of an arithmetic sequence can also be represented as the product of the average term in the sequence and the number of terms. Therefore, we have $60=(\text{average term of the sequence})$ $\cdot$ $(\text{number of terms})$. We know that the number of terms must be even, since if it was odd, then the sum of the sequence would be an odd number. We also know that the average term of the sequence must be even, because since there are an even number of terms, then the average term is the median of the two "middle" terms, which must be even. Therefore, we must have it such that $(\text{average term of the sequence})$ and $(\text{number of terms})$ multiply to 60, while both factors have to be even. This gives $4$ ordered pairs $(x,y)$ with $x$ being the average term of the sequence and $y$ being the number of terms: $(2, 30), (6, 10), (10, 6), (30, 2)$. However, $(2, 30)$ gives an arithmetic sequence with some numbers in the sequence being negative, and since the problem specifies we need positive integers, we discard this case. Similarly, $(6, 10)$ doesn't work either. The other cases work though, so we have a total of $\boxed{\textbf{(B) }2}$ cases. ~Baihly2024

 



\textbf{Solution 4 (Nice language, a little long)}

Let's write the possible pairs down, in the form $($\# of numbers, average of the sequence $)$. This will work, because the number of terms times the average is the sum of the numbers (proof provided at bottom, but it's intuitive). Now we have:
 \[(1,60), (2,30), (3,20), (4,15), (5,12), (6,10), (10,6), (12,5), (15,4), (20,3), (30,2), (60,1)\]

 That's a lot! Let's eliminate the obvious losers: We know $(1,60)$ is out, because we need at least two numbers. $(60, 1)$ is also out because that's $60$ odd integers that average to $1$, so some will have to be negative (and that violates the rules!). Checking out $(2,30)$, that works with $29$ and $31$. We can see $(10, 6)$ and anything to the right of that won't work, because they will have to contain more negative numbers to the left of them; $(10,6)$ has to have numbers $-3, -1, 1, 3, 5$ on the negative side. We want to make sure $(6, 10)$ hasn't turned to the negative side, so we try write it down: $5, 7, 9, 11, 13, 15$. That works!

 So far, we have: 
 \[\text{eliminated:} (1,60), (10,6), (12,5), (15,4), (20,3), (30,2), (60,1)\]
 \[\text{confirmed}: (2,30), (6,10)\]
 \[\text{checking}:(3,20), (4,15), (5,12)\]Wait a minute! Three odd consecutive integers can't average to an even! If we try: $((2x-1) + (2x+1) + (2x+3))/3 = (6x + 3)/3 = (2x+1)$, which is an odd number! That eliminates $(3,20)$ and $(5,12)$. Finally, we see that 4 odd numbers can't average to an odd, because $(2x-1) + (2x-3) + (2x+1) + (2x+3) = 8x$, an even number. So that eliminates the last one, and we are left with $2$ confirmed cases.

 \textbf{Proof used in solution 4 (trivial)} Say we have a list on integers $[x_1, x_2, x_3, ... , x_n]$. Then there are $n$ terms. The average of this list is $(x_1 + x_2 + ... + x_n)/n$, and when we multiply it by $n$, we get $x_1 + x_2 + ... + x_n$, the sum of the list. Therefore, the number of terms times the average of the terms, is the sum of the terms.

 


 



\textbf{Solution 5 (Difference of Squares)}

Recall that for integers $n \geq 0$, the sum of the first $n$ odd positive integers is $n^2$.  Therefore, the sum of the $a^{\textrm{th}}$ through $b^{\textrm{th}}$ odd numbers, where $a \leq b$ are positive integers, is $b^2-(a-1)^2$.

 Let $c = a-1$, so $c$ is an integer such that $0 < c < b$.  Then each series of consecutive odd positive integers summing to $60$ corresponds to a solution to the equation $b^2 - c^2 = 60$, where $b > c$ are positive integers.  Factoring yields $(b+c)(b-c) = 60$.  Since $b+c$ and $b-c$ must either be both even or both odd, they are both even.  That is, $b+c = 2m$ and $b-c = 2n$ for some  positive integers $m > n$.  Since $(2m)(2n) = 4mn = 60$, $mn = 15$.  The only values of $(m, n)$ satisfying this equation are $(15, 1)$ and $(5, 3)$.  The corresponding values of $(a, b)$ are $(15, 16)$ and $(3, 8)$, which yield $5 + 7 + 9 + 11 + 13 + 15$ and $29 + 31$, respectively.  Therefore, there are $\boxed{\textbf{(B) }2}$ such solutions.

 -j314andrews

 



\textbf{Solution 6 (Consecutive Number Addition)}

You know that adding consecutive numbers can be represented by $(x) + (x + 1) + (x + 2)...$. So, adding consecutive odd numbers would be represented by $(x) + (x + 2) + (x + 4)...$, since every other number is odd. However, every other number is also even. So you must be careful as to what number $x$ is so you know if this series will be even or odd. Split this problem up into $5$ cases (the answer choices really help out since the max is $5$). 

 Case 1 (two consecutive numbers): $(x) + (x + 2) = 60$, $2x+2 = 60$, $x=29$ which satisfies the odd and consecutive constraint.

 Case 2 (three consecutive numbers): $(x) + (x + 2) + (x + 4) = 60$, $3x+6 = 60$, $x=18$ which doesn't satisfy the odd and consecutive constraint.

 Case 3 (four consecutive numbers): $(x) + (x + 2) + (x + 4) + (x + 6) = 60$, $4x+12 = 60$, $x=12$ which doesn't satisfy the odd and consecutive constraint.

 Case 4 (five consecutive numbers): $(x) + (x + 2) + (x + 4) + (x + 6) + (x + 8) = 60$, $5x+20 = 60$, $x=8$ which doesn't satisfy the odd and consecutive constraint.

 Case 5 (six consecutive numbers): $(x) + (x + 2) + (x + 4) + (x + 6) + (x + 8) + (x + 10) = 60$, $6x+30 = 60$, $x=5$ which satisfies the odd and consecutive constraint.

 Therefore, there are 2 cases out of the 5 which satisfy the constraints, so there are $\boxed{\textbf{(B) }2}$ such solutions.

 -N.

 



\textbf{Solution 7 (straightforward)}

First, note that the sum of an odd number of odd integers is odd. Since $60$ is even, the number of consecutive odd integers, $n$, must be even.

 For $n=2$, the two numbers are centered around $\frac{60}{2}=30$, so they are $30\pm1$, giving $29$ and $31$. Their sum is $60$, so this works.

 For $n=4$, the numbers are centered around $\frac{60}{4}=15$, so they are $15\pm1$ and $15\pm3$, giving $12,14,16,18$. These are not odd, so this case is invalid.

 For $n=6$, the numbers are centered around $\frac{60}{6}=10$, so they are $10\pm1,\;10\pm3,\;10\pm5$, giving $5,7,9,11,13,15$. These are consecutive odd positive integers whose sum is $60$, so this works.

 For $n=8$, we have $\frac{60}{8}=7.5$, which is not an integer, so this case is impossible.

 For $n=10$, the numbers would be centered around $\frac{60}{10}=6$, giving $6\pm1,\;6\pm3,\;6\pm5,\;6\pm7,\;6\pm9$. This includes negative numbers, which are not allowed since the problem asks for positive integers.

 Thus, there are exactly $2$ valid representations, so the answer is $\boxed{\textbf{(B)}\ 2}$.

 



\textbf{Solution 8 (Factors and Mean)}

Let $a$ be the first odd positive integer, then 
 \[a + (a + 2) + (a + 4) + \dots + (a + 2(n-1)) = 60\] 
 \[na + (n-1)n = 60\]
 \[n^2 + (a-1) n = 60\]

 The mean of the consecutive odd numbers 
 \[n + (a-1) = \frac{60}{n}\]Note: $60=2^2 \times 3 \times 5$ ; $n$ is even as sum of odd number of odd numbers is odd ; $(a-1)$ is even

 The mean $\frac{60}{n} = n + (a-1)$ is an even integer and hence $n$ could not be a multiple of $4$.

 The possible values of $n = \{2, 2 \times 3, 2 \times 5, 2 \times 3 \times 5\}$, But when $n = 10, 30$ the mean $\frac{60}{n}$ will be smaller than $n$ such that $a$ could not be positive integer.

 Therefore, $n$ = $2, 6$ and the answer is $\boxed{\textbf{(B)}\ 2}$.

 


	\item (\textit{AMC 8 2026}) Miguel is walking with his dog, Luna. When they reach the entrance to a park, Miguel throws a ball straight ahead and continues to walk at a steady pace. Luna sprints toward the ball, which stops by a tree. As soon as the dog reaches the ball, she brings it back to Miguel. Luna runs $5$ times faster than Miguel walks. What fraction of the distance between the entrance and the tree does Miguel cover by the time Luna brings him the ball?


 
                    \begin{figure}[H]
                        \centering
                        \includegraphics[width=0.5\textwidth]{Problem/AMC_8_AMC_2026_19_Problem_Figure_1.png}
                    \end{figure}
                    


 $\textbf{(A) } \frac{1}{6}\qquad\textbf{(B) } \frac{1}{5}\qquad\textbf{(C) } \frac{1}{4}\qquad\textbf{(D) } \frac{1}{3}\qquad\textbf{(E) } \frac{2}{5}$


 
 \textbf{Answer: }D



\textbf{Solution}

Let Miguel's speed equal $s$. That makes Luna's speed $5s$. Let the distance between the park entrance and the tree be $d$. The total distance covered by Miguel and Luna is $s + 5s = 6s$, which is also equal to $2d$. Thus $6s = 2d \Rightarrow s = \boxed{ \textbf{(D) } \frac{1}{3}}d$.

 



\textbf{Solution 2}

The distance between the park entrance and the tree is $x$. Then, the rate that Miguel walks is $v$, and Luna runs at $5v$. As Luna runs the distance $x$ and back, she runs a total of $2x - m$ (where $m$ is the amount that Miguel walks). Then, because $d/s = t$ (where $d$ is distance, $t$ is time, and $s$ is speed), It takes Luna $\frac{2x - m}{5v}$ time to reach back to Miguel. Also, $d = st$, and by that time, Miguel traveled $m$, so $m = v\left(\frac{2x - m}{5v}\right)$, or $5m = 2x - m$, $6m = 2x$, $3m = x$. The total distance from the park to the tree is $3m$, and Miguel traversed $m$, giving us a total of $\frac{m}{3m} = \boxed{\frac{1}{3}}$ the distance traveled.

 



\textbf{Solution 3 (Set the Distance and Speed!)}

The distance doesn't seem to matter as if it did, then the problem would have told the distance, so lets say the distance is 30 units. Also the speed of Luna doesn't matter as the ratio of Luna's speed to Miguel's speed is 5:1, always the same. We can say that Luna goes 5 units a second, then that means that Miguel goes $\frac{5}{5} =1$ units a second. Luna will reach the tree after $\frac{30}{5} =6$ seconds. When Luna reaches the tree Miguel will be $6 \times 1 = 6$ units from the entrance of the park. We can set up the equations: $x+6$ and $-5x+30$, where x represents seconds as Miguel is 6 units from the entrance is walking away at a rate of 1 unit a second and Luna is 30 units from the entrance and is walking towards the entrance at a rate of 5 units a second. We can now set both equations equal together to find out how far away from the entrance they meet. $x+6=-5x+30 \implies 6x=24 \implies x=4$. This means that they meet 4 seconds after Luna reaches the tree, this means they meet after $4+6=10$ seconds as Luna takes 6 seconds to reach the tree. In $10$ seconds Miguel travels $10 \times 1=10$ units, and $\frac{10}{30} = \frac{1}{3}$, giving us $\boxed{\frac{1}{3}}$~ Shaarav14 for Carnage MS

 



\textbf{Solution 4}

Since Luna is 5 times faster than Miguel, Miguel gets $\frac{1}{5}$  of the way to the tree by the time Luna reaches the ball. The ratio of Miguel´s speed to Luna´s is 1:5, and the distance between Miguel and the tree is $\frac{4}{5}$ the distance between the entrance and the tree. So Miguel walks another $\frac{4}{5} \times \frac{1}{6}$ of the way from the entrance to the tree by the time Luna comes back to him. Therefore, Miguel walks $\frac{1}{5} + \frac{4}{30}$ =$\boxed{\frac{1}{3}}$ of the way to the tree from the entrance.-Pillowwillowtree01

 



\textbf{Solution 5 (Quick)}

Since Luna runs 5 times faster than Miguel walks, by the time Luna reaches the tree, Miguel has walked $\frac{1}{5}$ of the distance from the entrance to the tree. This eliminates answer choices $\text{(A)}$ and $\text{(B)}$.

 This is a multiple-choice problem with the remaining answers listed in increasing order, so we test the middle option, $\text{(D)}$, which corresponds to Miguel walking $\frac{1}{3}$ of the distance.

 If Miguel walks a total distance of $\frac{1}{3}$, then the additional distance he walks while Luna is running back toward him is $\frac{1}{3}-\frac{1}{5}=\frac{2}{15}$. Since Luna runs 5 times faster than Miguel walks, she covers $5\cdot\frac{2}{15}=\frac{2}{3}$ of the distance on her return.

 Thus, Luna runs the entire distance to the tree and $\frac{2}{3}$ of the distance back, while Miguel walks $\frac{1}{3}$ of the distance forward, so they meet at the same point. Therefore, the fraction of the distance Miguel covers is $\boxed{\frac{1}{3}}$.

 



\textbf{Solution 6}

Because Luna is running five times as fast as Miguel, when the two meet, Luna will have run five times as far. We can call the distance from the place where Miguel threw the ball to the tree $X$, and the distance that Miguel runs $Y$. When the two meet they will have run a combined total of $2X$. Miguel runs $Y$, and because she runs five times as fast, Luna runs $5Y$. There combined total is therefore $Y+5Y=6Y$. So $6Y=2x$. The value we are trying to find is $\frac{Y}{X}$

 We divide both sides of our equation by $6x$ to get $\frac{6Y}{6X}=\frac{2X}{6X}$. This can be simplified to get $\frac{Y}{X}=\frac{1}{3}$. So the fraction of the distance to the tree that Miguel runs is $\frac{Y}{X}=\boxed{\frac{1}{3}}$

 



\textbf{Solution 7 - Cheese 🧀}

If you do not know how to use normal math to solve this question, you can also use a ruler and utilize the provided diagram. This gives that the meeting point is $\frac{1}{3}$ in regard to distance to the tree, which means that the answer is $\boxed{ \textbf{(D) } \frac{1}{3}}$

 $\textbf{Note:}$ Even though test proctors, and essentially everybody, says that the diagrams and images in competition math competitions like AMC 8 are not necessarily drawn to scale, they almost always are. This question isn’t an exception. There are some exceptions, though (like the 2022 AMC 8 Problem 24.) So do not rely on this method, unless you have no other choice.

 


	\item (\textit{AMC 8 2026}) The land of Catania uses gold coins and silver coins. Gold coins are $1$ mm thick and silver coins are $3$ mm thick. In how many ways can Taylor make a stack of coins that is $8$ mm tall using any arrangement of gold and silver coins, assuming order matters? 


 $\textbf{(A) } 3\qquad\textbf{(B) } 7\qquad\textbf{(C) } 10\qquad\textbf{(D) } 13\qquad\textbf{(E) } 16$


 
 \textbf{Answer: }D



\textbf{Solution 1}

A stack containing $3$ or more silver coins is at least $9$ mm tall, so a stack of height $8$ mm may contain at most $2$ silver coins.

 An $8$ mm stack with $0$ silver coins contains $8$ gold coins.  There is $\binom{8}{0} = 1$ such stack.

 An $8$ mm stack with $1$ silver coin contains $5$ gold coins.  There are $\binom{6}{1} = 6$ such stacks.

 An $8$ mm stack with $2$ silver coins contains $2$ gold coins.  There are $\binom{4}{2} = 6$ such stacks.

 Thus there are $1+6+6=\boxed{\textbf{(D)}\ 13}$ different $8$ mm stacks.  

 



\textbf{Solution 2 (Recursion)}

Let $a_k$ be the number of different stacks of height $k$ mm.  Then $a_1 = 1$, as the only possible $1$ mm stack is a single gold coin.  Also $a_2 = 1$, as the only possible $2$ mm stack is two gold coins, while $a_3 = 2$, as the two possible $3$ mm stacks are either three gold coins or one silver coin.

 If $k \geq 4$, a stack of height $k$ mm can be formed by placing a stack of height $(k - 1)$ mm on a gold coin or placing a stack of height $(k - 3)$ mm on a silver coin.  Therefore, $a_k = a_{k-1} + a_{k-3}$.  So $a_4 = 3$, $a_5 = 4$, $a_6 = 6$, $a_7 = 9$, and $a_8 = \boxed{\textbf{(D)}\ 13}$.

 -j314andrews

 


	\item (\textit{AMC 8 2026}) Charlotte the spider is walking along a web shaped like a $5$-pointed star, shown in the figure below. The web has $5$ outer points and $5$ inner points. Each time Charlotte reaches a point, she randomly chooses a neighboring point and moves to that point. Charlotte starts at one of the outer points and makes $3$ moves (re-visiting points is allowed). What is the probability she is now at one of the outer points of the star?


 
                    \begin{figure}[H]
                        \centering
                        \includegraphics[width=0.5\textwidth]{Problem/AMC_8_AMC_2026_21_Problem_Figure_1.png}
                    \end{figure}
                    


 $\textbf{(A) } \frac{1}{5}\qquad\textbf{(B) } \frac{1}{4}\qquad\textbf{(C) } \frac{2}{5}\qquad\textbf{(D) } \frac{1}{2}\qquad\textbf{(E) } \frac{3}{5}$


 
 \textbf{Answer: }B



\textbf{Solution 1}

The graph undergoes \textbf{rotational symmetry} meaning when scaled on a factor $k$, here $k = 5$ of the points, it is the same as considering a fixed starting position of Charlotte. WLOG let this starting point be the topmost one. On the second turn, she cannot be on an outer point. Then she is forced to an inner point in a 1 chance. From any outer point, she must move to an outer point with a $\frac{1}{2}$ chance. Finally, from that outer point, she must move to an inner point with $\frac{1}{2}$ chance. Therefore, the probability is $\boxed{\frac{1}{4}}$.

 In general, the probabilistic depends on if the pentagram is on some $n$th moveset. For instance, $p_0 = 1$, and all $p_n = \frac{1}{2}(1-p_{n-1})$. The term $(1-p_{n-1})$ is of the point that to get to an outer point one must come from an inner point. 

 Note to SquishBJ from Pinotation: You are absolutely correct, thank you for taking the time to review my solution. Constructive counting is a swing and miss here due to the fact the question doesn't ask for strict tournaments rather cases based on dependency. Good catch :>

 



\textbf{Solution 2}

We begin by noticing that from an outer point, the next move must be to an inner point as both adjacent points are inner. Next, there are four adjacent points - two are inner, two are outer. Thus, there is a $\frac{1}{2}$ probability for both. The final move depends on which type of point was the last one - if it was outer, the next one would have to be inner, so that gives us a probability of $0$ for an outer point. However, if it was inner, there is, again, a $\frac{1}{2}$ probability that the next one is outer. Therefore our final probability is $1 \cdot (\frac{1}{2} \cdot \frac{1}{2} + \frac{1}{2} \cdot 0) = \frac{1}{2} \cdot \frac{1}{2} = \boxed{\textbf{(B) } \frac{1}{4}}$.

 



\textbf{Solution 3 (Markov Chains)}

Note that there are only two states due to five-fold symmetry. Hence, we can create a Markov Chain with two states: outer and inner points. Transitioning from an outer point to an inner point happens with probability $1$, and transitioning from an inner point to an outer point happens with probability $\frac{1}{2}$. The spider always changes from states, and it never stays on the same state in two consecutive moves. The spider starts on an outer point, so the probability at Move 0 for outer points is $1$.

 Enumerating the values for 3 moves using our recursive rules above:

 
 \[\begin{array}{|c|c|c|c|} \hline &&\\ \textbf{Move} & \textbf{Outer} & \textbf{Inner}\\ &&\\ \hline &&\\ 0 & 1 & 0\\ &&\\ \hline &&\\ 1 & 0 & 1\\ &&\\ \hline &&\\ 2 & \frac{1}{2} & \frac{1}{2}\\ &&\\ \hline &&\\ 3 & \frac{1}{4} & \frac{3}{4}\\ &&\\ \hline \end{array}\]We aim to find the probability of the spider being at an outer point on Move 3. Thus, the answer is $\boxed{\textbf{(B)} \frac{1}{4}}$.

 



\textbf{Solution 4 (Recursive)}

Let $p_n$ be the probability that Charlotte is at an outer point after $n$ moves.

 From the symmetry of the star graph we derive that the probabilities are:

 
                \begin{itemize}
                    \item Outer $\to$ Inner with probability $1$\item Inner $\to$ Outer with probability $\frac12$\item Inner $\to$ Inner with probability $\frac12$
                \end{itemize}


                If Charlotte is at an outer point after $n+1$ moves, she must have been at an inner point after $n$ moves and then moved outward. Therefore
 \[p_{n+1}=\frac12(1-p_n)\]

 The initial condition is
 \[p_0=1\]since she starts at an outer point.

 Rewrite the recurrence as
 \[p_{n+1}=\frac12-\frac12 p_n.\]

 


 First, we find the constant equilibrium solution $c$:
 \[c=\frac12-\frac12 c\]
 \[\frac32 c=\frac12\]
 \[c=\frac13.\]

 

 \[q_n=p_n-\frac13.\]

 Substituting into the recurrence,
 \[p_{n+1}=\frac12-\frac12\left(q_n+\frac13\right) =\frac13-\frac12 q_n.\]

 Thus
 \[q_{n+1}=-\frac12 q_n.\]

 We now have a clean equation.
 \[q_n=q_0\left(-\frac12\right)^n.\]

 Now we do some simple algebra.
 \[q_0=p_0-\frac13=1-\frac13=\frac23,\]
 \[q_n=\frac23\left(-\frac12\right)^n.\]
 \[p_n=q_n+\frac13=\frac13+\frac23\left(-\frac12\right)^n.\]

 Now we plug in $n=3$:
 \[p_3=\frac13+\frac23\left(-\frac12\right)^3 =\frac13-\frac1{12} =\frac14.\]

 Thus the probability Charlotte is at an outer point after $3$ moves is $\boxed{\frac14}$.

 \textbf{Remark} This solution turns a recursive sequence into a geometric sequence. It's overkill for this question, but it could be used if we are finding the probability after a large number of moves, e.g., 2026 moves.

 



\textbf{Solution 5 (Quick)}

Like solution 1 and 2 we know that we can choose one edge point to start on. From any edge point there are two possibilities to go to, we can just use one of these possibilities due to symmetry, making the problem easier to solve. Once we go to this point we can see that there is a $\frac12$ probability that you go to an edge point and a $\frac12$ probability that you stay in a center point. If you go to an edge point there is a 0 probability that you can end on an edge point. While if you stay in the middle, there is a $\frac12$ probability that you go to an edge point on your 3rd move. We can calculate this using symmetry without loss of generality making this problem much easier to solve. Finally our answer is $\frac12*\frac12=\frac14$ which is $\boxed{B:\frac14}$

 
(PS. Make sure to check whether you can use symmetry for a counting problem before you actually use it.)~SShreyas

 


	\item (\textit{AMC 8 2026}) The integers from $1$ to $25$ are arbitrarily separated into five groups of $5$ numbers each. The median of each group is identified. Let $M$ equal the median of the five medians. What is the least possible value of $M$?


 $\textbf{(A) } 9\qquad\textbf{(B) } 10\qquad\textbf{(C) } 12\qquad\textbf{(D) } 13\qquad\textbf{(E) } 14$


 
 \textbf{Answer: }A



\textbf{Solution 1}


                    \begin{figure}[H]
                        \centering
                        \includegraphics[width=0.5\textwidth]{Solution/AMC_8_AMC_2026_22_Solution_Figure_1.png}
                    \end{figure}
                    

 

 Consider using a five by five square table to solve this problem. The rows represent the five groups, the center column is the median of each group and the center square cell is the least possible value of \(M\). The only numbers that really matter is the ones in the shaded area, the rest can be ignored. In order to find the least possible value of \(M\) the smallest numbers, \(1 ... 9\), need to be in these nine square cells. \(9\) is the smallest number that is able to satisfy both median of a group and median of five medians. Final answer, the least possible value of \(M\) is $\boxed{\textbf{(A)}~9}.$ ~TutorJack

 



\textbf{Solution 2 (Quick and easy)}

If a group has median $m$, then we must have that $3$ of the numbers in that group are $\leq m$. Since there are 5 different groups, $3$ groups must have a median $\leq M$, so there are at least $3\cdot3=9$ numbers that are $\leq M$. Since there are at least $9$ numbers that are $\leq M$, we have $M$ at minimum $\boxed{\textbf{(A) }9}.$ ~Baihly2024

 



\textbf{Solution 3 (Solution 1 Rigorously)}

Let the groupings be $G_1, G_2, \dots, G_5$. Let the medians of each of the groupings be $m_1, m_2, m_3, m_4,$ and $m_5$ in that order. Let $M$ be the median of those medians. Then, we wish to minimize $M$. This can be possible if we have three elements in $G_1$ less than $M$, three elements in $G_2$ less than $M$, and three elements in $G_3$ less than an equal to $M$ are there (as the last number itself is $M$). Thus, the smallest way to choose $9$ such numbers is $\{1, 2, 3, 4, \dots, 9\}$, in which $M$ is the largest of these, and thus the answer is $\boxed{9}$.

 


 



\textbf{Solution 4}

If you want to prove a solution one of the easiest ways is just to list an example of it.A median has to be greater than the first two numbers of the group.Then we just list the groups out.This includes:

 \begin{align*}( 1,2,3,16,17) ,\ ( 4,5,6,18,19) ,\ ( 7,8,9,20,21) ,\ ( 10,11,12,22,23) ,\ ( 13,14,15,24,25) & & \end{align*}


 Then according to this triple we are able to find the $M$ is $\boxed{\textbf{(A) }9}$.

 


	\item (\textit{AMC 8 2026}) Lakshmi has $5$ round coins of diameter $4$ centimeters. She arranges the coins in $2$ rows on a table top, as shown below, and wraps an elastic band tightly around them. In centimeters, what will be the length of the band?


 
                    \begin{figure}[H]
                        \centering
                        \includegraphics[width=0.5\textwidth]{Problem/AMC_8_AMC_2026_23_Problem_Figure_1.png}
                    \end{figure}
                    


 $\textbf{(A)}\ 2\pi + 20 \qquad \textbf{(B)}\ \frac{5}{2}\pi + 20 \qquad \textbf{(C)}\ 4\pi + 20 \qquad \textbf{(D)}\ \frac{9}{2}\pi + 20 \qquad \textbf{(E)}\ 5\pi + 20$


 
 \textbf{Answer: }C



\textbf{Solution 1}

\begin{verbatim}
	unitsize(5mm);    draw(circle((2,2),2));    o));    draw(circle((10,2),2));    draw(circle((4,2+sqrt(12)),2));    draw(circle((8,2+sqrt(12)),2));    draw(arc((2,2),2,270,150) -- arc((4,2+sqrt(12)),2,150,90) .. arc((8,2+sqrt(12)),2,90,30) -- arc((10,2),2,30,-90) .. cycle, black+3bp);        fill((2,2) -- arc((2,2),2,270,150) -- cycle, mediumgray+opacity(0.4));    fill((4,2+sqrt(12)) -- arc((4,2+sqrt(12)),2,150,90) -- cycle, mediumgray+opacity(0.4));    fill((8,2+sqrt(12)) -- arc((8,2+sqrt(12)),2,90,30) -- cycle, mediumgray+opacity(0.4));    fill((10,2) -- arc((10,2),2,30,-90) -- cycle, mediumgray+opacity(0.4));    draw((2,2) -- (4,2+sqrt(12)) -- (8,2+sqrt(12)) -- (10,2) -- (6,2) -- cycle);    draw((2,2) -- arc((2,2),2,270,150) -- cycle);    draw((4,2+sqrt(12)) -- arc((4,2+sqrt(12)),2,150,90) -- cycle);    draw((8,2+sqrt(12)) -- arc((8,2+sqrt(12)),2,90,30) -- cycle);    draw((10,2) -- arc((10,2),2,30,-90) -- cycle);    draw((6,2) -- (6,0)); (Error making remote request. Unknown error\_msg)
\end{verbatim}



 In order to find the length of the elastic band, perpendicular lines are added to the figure to divide the band into sections for easier calculation. After the additions, the length of band consist of five identical lines and a full circle, the four shaded areas. The total length is five lines, the diameter of the coins, plus the circumference of a coin, so \(5d + 4\pi\). Finally, the length of the elastic band is $\boxed{\textbf{(C)}4\pi+20}.$~TutorJack

 



\textbf{Solution 2 (Pure Geometry)}

As shown in the figure above, we just connect some lines and do relations to find the straight parts of the bands (That's easy). Now, we need to find the four curves. The three central circles form a triangle with angle measures $60^\circ$. Extend the lines of that triangle and draw two perpendicular lines going through the top left circle's center. Then, it becomes obvious the angle measure for the two topmost arcs of the band is $60^\circ$. Thus, the length of the curve is $\frac{1}{6} \cdot \pi \cdot 2 \cdot 2 = \frac{2}{3}\pi$. There are two of these, giving us a total of $\frac{4}{3}\pi$. Now, the bottom two. We do the same process and this time find out the measure of the arcs is $120^\circ$, meaning the length of the curve is $\frac{1}{3} \cdot 2 \cdot \pi \cdot 2 = \frac{4}{3} \pi$. There are two such curves, and thus the total length is $\frac{8}{3} \pi$. The total curvature length is $\frac{12}{3} \pi = 4\pi$. The straight parts is just additive to $4+8+4+4 = 20$. Thus, the length of the band is $\boxed{20+4\pi}$.

 



\textbf{Solution 3 (Using the answer choices, quickest and best for competitions)}

If you look at the answer choices you will notice that all the constants are the same, meaning that we just have to find the coefficient of pi. To do this we can form an isosceles trapezoid with the centers of the four corner circles.

 
                    \begin{figure}[H]
                        \centering
                        \includegraphics[width=0.5\textwidth]{Solution/AMC_8_AMC_2026_23_Solution_Figure_1.png}
                    \end{figure}
                    

 

 
Once we do this we can split the trapezoid into three equilateral triangles, after doing this we can solve for the curved lengths. The bottom 2 curved lengths are each $120^\circ$ as $360-90-90-60=120,$ because this is an isosceles trapezoid the top two angles must be the same, thus they are both $60^\circ$ as $\frac{360-(120+120)}{2}=60$ (We can also solve for this like we did for the bottom two by doing $360-90-90-60-60=60$) $60+60+120+120=360$ so the total curved part is the same as the circumference of the circle which is diameter times $pi$ which is 4*pi which = $4\pi,$ so the answer is $\boxed{\textbf{(C)}4\pi+20}.$.

 PS. For problems like these where a band is wrapped around congruent circles, the curved part is always equal to the circumference of a circle. 

 



\textbf{Solution 5 (better visualization)}


                    \begin{figure}[H]
                        \centering
                        \includegraphics[width=0.5\textwidth]{Solution/AMC_8_AMC_2026_23_Solution_Figure_2.png}
                    \end{figure}
                    

 There are $4$ fourths of a circle with radius $2$, so there is one circle with radius $2$. Finding the circumference gives $4\pi$, and adding the $4$'s created by matching the centers of one circle to another by using the radius property gives (it's always equal if it comes from the center) $\boxed{\textbf{(C)}4\pi+20}.$

 -N.

 



\textbf{Solution 6 (Cheese🧀🧀🧀🧀🧀)}

We know that the coefficient of $\pi$ must be $2 \times 2$ because the angles in the trapezoid add to $360^\circ$, the circles are the same size, and the formula for the circumference of a circle is $\pi d$. The straight bits add to the rational bit and the curved bits add to the irrationals (with $\pi$). Therefore, the coefficient is $4\pi$ so the correct answer is $\boxed{\textbf{(C)}4\pi+20}.$

 


	\item (\textit{AMC 8 2026}) The factorial notation $n!$ is defined as the product of the first $n$ positive integers. (For example, $3! = 1 \cdot 2 \cdot 3 = 6$.) Define the superfactorial of a positive number, denoted by $n^{!}$, to be the product of the factorials of the first $n$ integers. (For example, $3^{!} = 1! \cdot 2! \cdot 3! = 12$.) How many factors of $7$ appear in the prime factorization of $51^{!}$, the superfactorial of $51$?   


 $\textbf{(A)}\ 147 \qquad \textbf{(B)}\ 150 \qquad \textbf{(C)}\ 156 \qquad \textbf{(D)}\ 168 \qquad \textbf{(E)}\ 171$


 
 \textbf{Answer: }E



\textbf{Solution 1 (Solution 2 simplified) - Come and Go Solution}

From $7! - 13!$, there is one power of $7$. From $14! - 20!$, there are $2$ powers. $21! - 27!$, $3$ powers. $28! - 34!$, $4$ powers. From $35! - 41!$, $5$ powers. From $42! - 48!$, $6$ powers. Finally, from $49!$ to $51!$, there are $8$ powers (as each one contributes $7^2$). Thus, in a superfactorial, the number of powers is just $7+14+21+28+35+42+24$ powers total, or $\boxed{171}$.

 In another way, we can think of it like 7 multiples. For example, from seven to thirteen there happens to exist 7 7's. Later from 14 to 20 there happens to exist 7*2 sevens. When we continue to move on to 49 and 51, there exists an extra seven because of the squares. Adding them up we should get 171 as our answer.

 Note for Jellycat88: That logic is flawed. Definition of multiples always implies one factor of seven. your idea works as long as the value is not a perfect square. for instance, $14!$ through $21!$, each corresponds two multiples of seven, those being 7 and 14. when you count that extra square, you are counting a power not a multiple, because from 49 to 51, we obtain only seven multiples, 7,14,21,28,35,42,49. However, from 49 through 51, we obtain eight \textbf{powers}. it is a minor yet huge difference, which is why your logic is slightly flawed.

 



\textbf{Solution 2 - Thought Process}

We begin by noticing that for the first few factorials ($1!$ to $6!$) there are no factors of $7$, because the first multiple of $7$ is $7$ itself. So, these numbers contribute $0$ sevens.

 The next $7$ factorials ($7!$ to $13!$) all have one factor of $7$ in them, since every number less than or equal to $n$ is multiplied in $n!$. These all contribute $7$ sevens.

 For the next $7$ factorials ($14!$ to $20!$), there are two factors of $7$ (namely $7$ and $14$). These numbers contribute $14$ sevens.

 If we continue with this method, we see that once we hit $7k!$ we have $7k$ sevens contributed, since there are $7$ numbers in between. This will work from $1!$ to $48!$, because when we get to $49! = (7^2)!$, we have one extra $7$ for every number, since we have an extra factor of $7$.

 This is only for $3$ numbers, $49!$, $50!$, and $51!$ (our final factorial). They have $7+1 = 8$ factors of $7$, because they have $7$ multiples of $7$ and one extra from $49$. These three contribute a total of $3 \cdot 8 = 24$ sevens.

 Finally, sum these numbers. Most are just the first $6$ multiples of $7$, and we can calculate this sum to be $0 + 7 + 14 + 21 + \dots + 42 = 7 \cdot \frac{6(6+1)}{2} = 7 \cdot \frac{42}{2} = 7 \cdot 21 = 147$. The final step is to add the final $3\cdot 8 = 24$ sevens to get our answer, $147 + 24 = \boxed{ \textbf{(E) } 171}$.

 



\textbf{Solution 3}

If $51^! = 1! \cdot 2! \cdot ... \cdot 51! = (1) (1 \cdot 2)(1 \cdot 2 \cdot 3)(...)(1 \cdot 2 \cdot 3 \cdot ... \cdot 51)$ is expanded as such, each integer $k \leq 51$ will appear once in each of the expansions of $\{k!, ..., 51!\}$, for a total of $52-k$ appearances.  

 Therefore, in this expansion, $7$ appears $45$ times, $14$ appears $38$ times, $21$ appears $31$ times, $28$ appears $24$ times, $35$ appears $17$ times, $42$ appears $10$ times, and $49$ appears $3$ times.  Except for $49$, which contributes $2$ copies of the factor $7$, each of these terms contributes $1$ copy of the factor $7$, for a total of $45 + 38 + 31 + 24 + 17 + 10 + 3 \cdot 2 = \boxed{ \textbf{(E) } 171}$ copies of the factor $7$.

 
-j314andrews

 


	\item (\textit{AMC 8 2026}) In an equiangular hexagon, all interior angles measure $120^\circ$. An example of such a hexagon with side lengths $2$, $3$, $1$, $3$, $2$, and $2$ is shown below, inscribed in equilateral triangle $ABC$. Consider all equiangular hexagons with positive integer side lengths that can be inscribed in $\triangle ABC$, with all six vertices on the sides of the triangle. What is the total number of such hexagons? Hexagons that differ only by a rotation or a reflection are considered the same.


 
                    \begin{figure}[H]
                        \centering
                        \includegraphics[width=0.5\textwidth]{Problem/AMC_8_AMC_2026_25_Problem_Figure_1.png}
                    \end{figure}
                    


 $\textbf{(A) } 4\qquad\textbf{(B) } 5\qquad\textbf{(C) } 6\qquad\textbf{(D) } 7\qquad\textbf{(E) } 8$


 
 \textbf{Answer: }E



\textbf{Solution 1 (Recursion)}

Let the side lengths of the hexagon be $a, b, c, d, e, f$. Obviously, each one of these is to be in between 1, 2, or 3 (see solution 2 for proof). Every 2 can be followed by a 2 or a 3. Every 3 must be followed by a 1 or 2, and every 1 must be followed by a 3. The question is now how many strings of length 6 satisfy these conditions, in which the first and last numbers also satisfy these conditions.

 To count that, we use recursion. Let $a_n$ be a string of length $n$ ending in 1, $b_n$ is a string of length $n$ ending in 2, and $c_n$ is a string that ends in 3. For $a_{n+1}$, we have that it must come from $c_n$, so $a_{n+1} = c_n$. Additionally, $b_{n+1} = b_n + c_n$, and $c_{n+1} = a_n$. Then, $a_1 = 1$, $b_1 = 1$, and $c_1 = 1$. Thus, $a_2 = 1$, $b_2 = 2$, and $c_2 = 1$; $a_3 = 1$, $b_3 = 3$, and $c_3 = 1$; $a_4 = 1$, $b_4 = 4$, $c_4 = 1$; $a_5 = 1$, $b_5 = 5$, $c_5 = 1$; and finally $a_6 = 1$, $b_6 = 6$, and $c_6 = 1$.

 Our answer is $a_6 + b_6 + c_6$, or $1+6+1 = \boxed{8}$.

 



\textbf{Solution 2}

An equiangular hexagon has all interior angles $120^\circ$, so consecutive sides turn by $60^\circ$ each time.Place an equilateral triangle $ABC$ of side length $6$ (this is the side length in the given example since$1+3+2=6$ along each side).

 Since all six vertices of the hexagon lie on the sides of $\triangle ABC$, each side of $\triangle ABC$contains exactly two consecutive vertices of the hexagon. Hence three alternating sides of the hexagon liealong the sides of $\triangle ABC$.

 Look near vertex $A$.Between the two sides $AB,AC$ and the hexagon there is a triangular gap whose vertex is $A$.Because $\angle A=60^\circ$ and the hexagon makes $120^\circ$ interior angles, each of the two angles of the gapadjacent to the hexagon is $60^\circ$. Therefore the gap at $A$ is an equilateral triangle.Do the same at $B$ and $C$: the hexagon is obtained by ``cutting off\textit{ three equilateral corner triangles.}

 Let the side lengths of the three cut-off equilateral triangles at $A,B,C$ be
 \[x,\ y,\ z \quad (\text{positive integers}).\]Along side $AB$ of length $6$, the two cuts remove lengths $x$ (near $A$) and $y$ (near $B$),so the remaining middle segment on $AB$ (which is a side of the hexagon) has length $6-x-y$.Similarly the other two middle segments on $BC$ and $CA$ have lengths $6-y-z$ and $6-z-x$.

 Thus, going around the hexagon, its six side lengths are
 \[(6-x-y,\ y,\ 6-y-z,\ z,\ 6-z-x,\ x),\]so all are positive if
 \[x+y\le 5,\qquad y+z\le 5,\qquad z+x\le 5,\]with $x,y,z\in \mathbb{Z}_{>0}$.

 Two such hexagons that differ by a rotation or reflection of the equilateral triangle correspond exactlyto permuting $(x,y,z)$, so we count solutions up to permutation. We may assume
 \[1\le x\le y\le z.\]

 Now enumerate all triples $(x,y,z)$ with $x\le y\le z$ satisfying the pair-sum constraints.

 If $x=1$, then $1+y\le 5\Rightarrow y\le 4$ and $1+z\le 5\Rightarrow z\le 4$.Also $y+z\le 5$.The valid triples are:
 \[(1,1,1),\ (1,1,2),\ (1,1,3),\ (1,1,4),\ (1,2,2),\ (1,2,3).\](For example, $(1,2,4)$ fails since $2+4>5$.)

 If $x=2$, then $2+y\le 5\Rightarrow y\le 3$ and $2+z\le 5\Rightarrow z\le 3$.Also $y\le z$ and $y+z\le 5$.The valid triples are:
 \[(2,2,2),\ (2,2,3).\]

 If $x\ge 3$, then $x+y\ge 3+3=6>5$, impossible.

 So there are
 \[6+2=8\]unordered triples, hence $8$ distinct hexagons up to rotation/reflection.

 $\boxed{8}$

 


 



\textbf{Solution 3 (Casework)}

We can already get $(1,1,1)$ and $(2,2,2)$ if we want equal sides (the case $(3,3,3)$ doesn't work as it makes a triangle)

 Then, we can compute $(1,1,2)$, $(2,2,3)$, $(2,2,1)$, $(1,1,4)$, $(1,1,3)$ for cases of the same side length (and also two of the same)

 Finally, we find the cause $(1,2,3)$ for all unique numbers.

 This gives us a total of $\fbox{(E)\ 8}$ total cases.

 



\textbf{Solution 4 (Burnside's Lemma, unnecessarily long)}

Let $X$ be the set of all possible orientations of the hexagon (where reflections and rotations are different), let $\mathcal{O}$ denote the set of all orbits of $X$ under $D_3$ where $D_3$ is the dihedral group of order $3$. Burnside's tells us that 
 \[|\mathcal{O}| = \frac{1}{|D_3|}\sum_{g \in D_3} |\text{Fix}(g)|\] where $\text{Fix}(g) = \{x \mid g \cdot x = x\}$, the fixed points under the action $g$. $|\mathcal{O}|$ is the solution to our problem in this case, as every rotation or reflection of the hexagon is represented in $D_3$. More accurately, $D_3$ represents the symmetries of the triangle, and the hexagon is the triangle with the vertices cut out.

 First, let's find $|X|$. Letting the sides of the hexagon be $a,b,c,d,e,f$, we note that $b=6-a-c$,$d=6-c-e$,$f=6-e-a$. Since $0<a,b,c,d,e,f<6$, $1\leq a+c,c+e,e+a\leq 5$. We can now do casework:

 $a=1$:We have $1 \leq c \leq 4$, $1 \leq e \leq 4$, and $1 \leq c+e \leq 5$. Letting $x=c$ and $y=e$, we manually check the lattice points in a square in the $xy$-plane with coordinates $(1,1)$ to $(4,4)$ for $x+y > 5$. $6$ elements satisfy this inequality, so they are the solutions $(c,e)$ which don't work - therefore, there are $16-6=10$ working solutions $(c,e)$ satisfying our inequalities.

 $a=2$:$1 \leq c \leq 3$, $1 \leq e \leq 3$, $1 \leq c+e \leq 5$. In a similar fashion, in the $3 \times 3$ square of lattice points only $(3,3)$ fails to meet $c+e\leq 5$ so $8$ solutions work from here.

 $a=3$:$1 \leq c \leq 2$, $1 \leq e \leq 2$, $1 \leq c+e \leq 5$. Obviously every solution works here, so $4$ solutions add to the count.

 $a=4$:Here, $(c,e)=(1,1)$ is the only working solution.

 The final count is $23$. Therefore, $|X|=23$.

 
There are three conjugacy classes of $D_3$: $\{e\},\{R,R^{-1}\},$ and $\{T,TR,TR^2\}$. We picture $D_3$ as a triangle with each of the vertices labeled $a,c,e$ in some order. (The hexagon is the equilateral triangle of side 6 with smaller equilateral triangles of length $a,c,e$ cut out of the corners. $\text{Fix}(g)=\text{Fix}(h)$ if $g$ and $h$ are in the same conjugacy class, so it suffices to find $|\text{Fix}(e)|$, $|\text{Fix}(R)|$, and $|\text{Fix}(T)|$. 

 Every "vertex" is in its own orbit after the $e$-action, meaning that every element is completely fixed, so $|\text{Fix}(e)=|X|$. 

 The orbit of $R$ contains all the vertices, meaning for an element to be fixed, every vertex must be the same, resulting in two triangles: $(1,1,1)$ and $(2,2,2)$. 

 Finally, the $T$-action fixes one of the vertices and swaps the other two. WLOG let $a$ be the fixed vertex and ${c,e}$ be swapped - meaning that they must be equal for the triangle to be fixed. Our equations become $1 \leq a+c \leq 5$, $1 \leq 2c \leq 5$, or that $c \in {1,2}$. $c=1$ generates $4$ solutions, while $c=2$ generates $3$. Therefore, $7$ triangles are fixed by the $T$-action, and $|\text{Fix}(T)|=7$.

 Plugging all this into Burnside's, \begin{align*}|\mathcal{O}| &= \frac{1}{|D_3|}\sum_{g \in D_3} |\text{Fix}(g)| \\|\mathcal{O}| &= \frac{1}{6}(1 \cdot 23+2 \cdot 2 + 3 \cdot 7) \\|\mathcal{O}| &= \frac{1}{6}(48) \\|\mathcal{O}| &= \fbox{(E) 8}.\end{align*}


 



\textbf{Solution 5 (Intuitive, Simple)}

Observe that the side length of $\triangle ABC$ is $6$. A hexagon is equiangular if and only if all ``outer triangles,\textit{ formed by two adjacent vertices of the }hexagon and one vertex of the equilateral triangle, are equilateral. We perform casework on each equilateral triangle and obtain $8$ different configurations. Hence $\boxed{8}$.

 


\end{enumerate}